% ============================================================
%  Chapitre 6 — Logique de Description
% ============================================================
\chapter{Logique de Description}

\section{Fondements théoriques}

\subsection{Positionnement}

La \textbf{logique de description} (\textsc{dl}) est une famille de formalismes
de représentation des connaissances dédiés à la modélisation de domaines
structurés en \emph{concepts} (classes) et \emph{rôles} (propriétés).
Elle constitue la base théorique des ontologies OWL~2.

\subsection{Vocabulaire et constructeurs}

\begin{table}[H]
\centering
\caption{Constructeurs de la logique de description $\mathcal{SROIQ}$ (OWL 2)}
\begin{tabular}{lll}
\toprule
\textbf{Constructeur} & \textbf{Notation} & \textbf{Signification} \\
\midrule
Intersection   & $C \dlinter D$         & appartient à $C$ et à $D$ \\
Union          & $C \dlunion D$         & appartient à $C$ ou à $D$ \\
Complément     & $\dlcomp{C}$           & n'appartient pas à $C$ \\
Restriction ex.& $\dlsome{R}{C}$        & a au moins un $R$-successeur dans $C$ \\
Restriction un.& $\dlall{R}{C}$         & tous les $R$-successeurs sont dans $C$ \\
Concept vide   & $\bot$                 & concept insatisfaisable \\
Concept universel & $\top$             & tout individu \\
\bottomrule
\end{tabular}
\end{table}

\subsection{TBox et ABox}

Une ontologie se compose de deux composantes :

\begin{description}
  \item[TBox (terminologique)] — ensemble d'axiomes définissant les concepts et
    leurs relations.
    Un \emph{axiome d'équivalence} $C \equiv D$ (définition complète) ou
    un \emph{axiome de subsomption} $C \sq D$ (condition nécessaire).

  \item[ABox (assertionnelle)] — ensemble d'assertions sur les individus :
    $C(a)$ (l'individu $a$ appartient au concept $C$) et
    $R(a,b)$ ($a$ est lié à $b$ par le rôle $R$).
\end{description}

\subsection{Hypothèse du monde ouvert (OWA)}

Contrairement à la logique des défauts (monde fermé implicite), la logique de
description adopte l'\textbf{hypothèse du monde ouvert} (\textsc{owa}) :
ce qui n'est pas asserté n'est ni vrai ni faux — il pourrait exister des
informations non encore connues.

\begin{tcolorbox}[notebox, title={OWA en pratique}]
Si l'ABox ne contient pas \code{signe\_chez(flenn, x)}, le raisonneur OWL ne
peut pas conclure que flenn n'a aucun label. Il faut soit une assertion
négative explicite, soit un axiome de complément, pour forcer cette conclusion.
Dans ce projet, on supplée par une \textbf{assertion explicite de classe}.
\end{tcolorbox}

\subsection{Raisonnement — HermiT}

Le raisonneur \textbf{HermiT} implémente un algorithme de tableaux pour OWL~2
(\textsc{sroiq}).
Il supporte les services de raisonnement suivants :

\begin{itemize}
  \item \code{isEntailed(axiom)} — vérifier si un axiome est conséquence logique
        de l'ontologie ;
  \item \code{getInstances(C, direct)} — obtenir les individus de la classe $C$ ;
  \item \code{getSubClasses(C, direct)} — obtenir les sous-classes de $C$ ;
  \item \code{isConsistent()} — vérifier la cohérence globale de l'ontologie.
\end{itemize}

\section{Application — Industrie musicale}

\subsection{TBox}

\begin{tcolorbox}[defbox, title={TBox musicale (sélection)}]
\begin{align*}
\text{Artiste}           &\equiv \text{Personne} \dlinter \dlsome{\text{produit}}{\text{Oeuvre}} \\
\text{Rappeur}           &\equiv \text{Artiste} \dlinter \dlsome{\text{pratique}}{\text{Rap}}
                                  \dlinter \dlall{\text{produit}}{\text{TitreRap}} \\
\text{RappeurAlgerien}   &\equiv \text{Rappeur} \dlinter \dlsome{\text{chante\_en}}{\text{Darija}} \\
\text{Label}             &\equiv \text{Organisation} \dlinter \dlsome{\text{distribue}}{\text{Album}} \\
\text{MajorLabel}        &\equiv \text{Label} \dlinter \dlsome{\text{signe}}{\text{PlusDe100Artistes}} \\
\text{IndependantLabel}  &\equiv \text{Label} \dlinter \dlcomp{\text{MajorLabel}} \\
\text{AlbumCertifie}     &\equiv \text{Album} \dlinter \dlsome{\text{reçoit}}{\text{Certification}} \\
\text{AlbumOr}           &\equiv \text{AlbumCertifie} \dlinter \dlsome{\text{reçoit}}{\text{CertificationOr}} \\
\text{ArtisteSign\'{e}} &\equiv \text{Artiste} \dlinter \dlsome{\text{signe\_chez}}{\text{Label}} \\
\text{ArtisteMajor}      &\equiv \text{Artiste} \dlinter \dlsome{\text{signe\_chez}}{\text{MajorLabel}} \\
\text{ArtisteIndependant}&\equiv \text{Artiste} \dlinter \dlcomp{\dlsome{\text{signe\_chez}}{\text{MajorLabel}}} \\
\text{ArtisteLibreDroits}&\equiv \text{Artiste} \dlinter \dlcomp{\dlsome{\text{signe\_chez}}{\text{Label}}}
\end{align*}
\end{tcolorbox}

\subsection{Hiérarchie de concepts}

\begin{figure}[H]
\centering
\begin{tikzpicture}[
  node distance=1.4cm and 2.2cm,
  every node/.style={concept}
]
  \node (Personne)  {Personne};
  \node (Artiste)   [below=of Personne]  {Artiste};
  \node (Rappeur)   [below left=of Artiste]  {Rappeur};
  \node (Chanteur)  [below right=of Artiste] {Chanteur};
  \node (RappAlg)   [below=of Rappeur]   {RappeurAlgerien};
  \node (ArtMajor)  [right=3cm of Rappeur]  {ArtisteMajor};
  \node (ArtSigne)  [above=of ArtMajor]  {ArtisteSigne};
  \node (ArtIndep)  [right=of ArtMajor]  {ArtisteIndependant};
  \node (ArtLibre)  [right=of ArtIndep]  {ArtisteLibreDroits};

  \draw[arrow] (Artiste) -- (Personne);
  \draw[arrow] (Rappeur) -- (Artiste);
  \draw[arrow] (Chanteur) -- (Artiste);
  \draw[arrow] (RappAlg) -- (Rappeur);
  \draw[arrow] (ArtMajor) -- (ArtSigne);
  \draw[arrow] (ArtSigne) -- (Artiste);
  \draw[arrow] (ArtIndep) -- (Artiste);
  \draw[arrow] (ArtLibre) -- (Artiste);
\end{tikzpicture}
\caption{Hiérarchie TBox — domaine musical ($A \to B$ signifie $A \sq B$)}
\label{fig:tbox-music}
\end{figure}

\subsection{ABox}

\begin{center}
\begin{tabular}{ll}
\toprule
\textbf{Assertion de classe} & \textbf{Assertions de propriétés} \\
\midrule
$\text{RappeurAlgerien}(\text{soolking})$ & $\text{signe\_chez}(\text{soolking}, \text{universal})$ \\
$\text{RappeurAlgerien}(\text{feu})$      & $\text{signe\_chez}(\text{feu}, \text{believe})$ \\
$\text{RappeurAlgerien}(\text{tif})$      & $\text{produit}(\text{soolking}, \text{guerilla})$ \\
$\text{RappeurAlgerien}(\text{flenn})$    & $\text{recoit}(\text{guerilla}, \text{cert\_or})$ \\
$\text{MajorLabel}(\text{universal})$     & $\text{ArtisteLibreDroits}(\text{tif})$ \\
$\text{IndependantLabel}(\text{believe})$ & $\text{ArtisteLibreDroits}(\text{flenn})$ \\
$\text{Album}(\text{guerilla})$           & \\
$\text{CertificationOr}(\text{cert\_or})$ & \\
\bottomrule
\end{tabular}
\end{center}

\subsection{Requêtes et résultats}

\begin{table}[H]
\centering
\caption{Résultats des requêtes HermiT — domaine musical}
\begin{tabularx}{\linewidth}{clcX}
\toprule
\textbf{\#} & \textbf{Requête} & \textbf{Résultat} & \textbf{Justification} \\
\midrule
R1 & $\text{ArtisteIndependant}(\text{soolking})$
   & \textcolor{red}{\textbf{FAUX}}
   & soolking $\in \dlsome{\text{signe\_chez}}{\text{MajorLabel}}$,
     contradiction avec $\dlcomp{}$ \\
R2 & $\text{ArtisteLibreDroits}(\text{feu})$
   & \textcolor{red}{\textbf{FAUX}}
   & feu $\in \dlsome{\text{signe\_chez}}{\text{IndependantLabel} \sq \text{Label}}$ \\
R3 & $\text{ArtisteLibreDroits}(\text{tif})$
   & \textcolor{codegreen}{\textbf{VRAI}}
   & asserté en ABox \\
R4 & $\text{ArtisteLibreDroits}(\text{flenn})$
   & \textcolor{codegreen}{\textbf{VRAI}}
   & asserté en ABox \\
R5 & $\text{AlbumOr}(\text{guerilla})$
   & \textcolor{codegreen}{\textbf{VRAI}}
   & recoit(guerilla, certOr), certOr $\sq$ Certification \\
R6 & $\text{instances}(\text{ArtisteSigne})$
   & \{soolking, feu\}
   & signe\_chez Label pour les deux \\
R7 & $\text{instances}(\text{ArtisteLibreDroits})$
   & \{tif, flenn\}
   & assertés, pas de signe\_chez \\
R8 & $\text{ArtisteMajor} \sq \text{ArtisteSigne}$
   & \textcolor{codegreen}{\textbf{VRAI}}
   & MajorLabel $\sq$ Label $\Rightarrow$ subsomption \\
R9 & Disjoint(ArtisteIndependant, ArtisteMajor)
   & \textcolor{codegreen}{\textbf{VRAI}}
   & $\dlcomp{\dlsome{R}{M}} \dlinter \dlsome{R}{M} = \bot$ \\
\bottomrule
\end{tabularx}
\end{table}

\subsection{Démonstration de non-monotonie OWL}

\begin{tcolorbox}[resultbox, title={R10 — Non-monotonie pour flenn}]
\textbf{Ontologie AVANT :} flenn : ArtisteLibreDroits (asserté) \\
\quad HermiT : \code{isEntailed(ArtisteLibreDroits(flenn))} $=$ \textbf{VRAI}

\medskip
\textbf{Ontologie APRÈS :}
\begin{enumerate}
  \item Retrait de l'assertion \code{ArtisteLibreDroits(flenn)}
  \item Ajout de \code{signe\_chez(flenn, universal)}
\end{enumerate}
\quad universal : MajorLabel $\sq$ Label
$\Rightarrow$ flenn $\in \dlsome{\text{signe\_chez}}{\text{Label}}$
$\Rightarrow$ flenn $\notin$ ArtisteLibreDroits \\
\quad HermiT : \code{isEntailed(ArtisteLibreDroits(flenn))} $=$ \textbf{FAUX}

\medskip
\textbf{Lien avec DefaultLogicMusic :}
la nouvelle assertion \code{signe\_chez(flenn, universal)}
joue le rôle du fait \code{SigneLabel(flenn)} qui déclenche le défaut $d_2$
et bloque $d_1$, invalidant \code{LibreDroits(flenn)}.
\end{tcolorbox}

\section{Application — Système judiciaire}

\subsection{TBox}

\begin{tcolorbox}[defbox, title={TBox judiciaire (sélection)}]
\begin{align*}
\text{Magistrat}           &\equiv \text{ActeurJudiciaire} \dlinter \dlsome{\text{siege\_a}}{\text{Tribunal}} \\
\text{Juge}                &\equiv \text{Magistrat} \dlinter \dlsome{\text{preside}}{\text{Audience}} \\
\text{Procureur}           &\equiv \text{Magistrat} \dlinter \dlsome{\text{requiert}}{\text{Peine}} \\
\text{AccuseMineur}        &\equiv \text{Accuse} \dlinter \dlsome{\text{a\_age}}{\text{MoinsDe18Ans}} \\
\text{AccuseRecidiviste}   &\equiv \text{Accuse} \dlinter \dlsome{\text{a\_antecedent}}{\text{Condamnation}} \\
\text{AccusePresumInnocent}&\equiv \text{Accuse} \dlinter \dlcomp{\text{Condamne}} \\
\text{AccuseAvecTemoignage}&\equiv \text{Accuse} \dlinter \dlsome{\text{beneficie}}{\text{TemoignageFav.}} \\
\text{CasSimple}           &\equiv \text{Accuse} \dlinter \dlcomp{\text{AccuseRecidiviste}}
                                   \dlinter \dlcomp{\text{AccuseMineur}} \\
\text{CasComplexe}         &\equiv \text{Accuse} \dlinter
                                   (\text{AccuseRecidiviste} \dlunion \text{AccuseMineur})
\end{align*}
\end{tcolorbox}

\begin{figure}[H]
\centering
\begin{tikzpicture}[node distance=1.3cm and 2cm, every node/.style={concept}]
  \node (Personne)   {Personne};
  \node (AJ)         [below=of Personne]         {ActeurJudiciaire};
  \node (Accuse)     [below left=of AJ]           {Accusé};
  \node (Magistrat)  [below right=of AJ]          {Magistrat};
  \node (Juge)       [below left=of Magistrat]    {Juge};
  \node (Proc)       [below right=of Magistrat]   {Procureur};
  \node (Mineur)     [below left=of Accuse]       {AccuséMineur};
  \node (Recid)      [below=of Accuse]            {AccuséRécid.};
  \node (Innocent)   [below right=of Accuse]      {AccuséPrésumInnocent};
  \node (Simple)     [below=of Mineur]            {CasSimple};
  \node (Complexe)   [below=of Recid]             {CasComplexe};

  \draw[arrow] (AJ) -- (Personne);
  \draw[arrow] (Accuse) -- (Personne);
  \draw[arrow] (Magistrat) -- (AJ);
  \draw[arrow] (Juge) -- (Magistrat);
  \draw[arrow] (Proc) -- (Magistrat);
  \draw[arrow] (Mineur) -- (Accuse);
  \draw[arrow] (Recid) -- (Accuse);
  \draw[arrow] (Innocent) -- (Accuse);
  \draw[arrow] (Simple) -- (Accuse);
  \draw[arrow] (Complexe) -- (Accuse);
\end{tikzpicture}
\caption{Hiérarchie TBox — domaine judiciaire}
\label{fig:tbox-justice}
\end{figure}

\subsection{Requêtes et résultats}

\begin{table}[H]
\centering
\caption{Résultats des requêtes HermiT — domaine judiciaire}
\begin{tabularx}{\linewidth}{clcX}
\toprule
\textbf{\#} & \textbf{Requête} & \textbf{Résultat} & \textbf{Justification} \\
\midrule
R1 & AccusePresumInnocent(ali)    & \textcolor{codegreen}{\textbf{VRAI}}  & asserté en ABox \\
R2 & AccuseAvecTemoignage(ali)    & \textcolor{codegreen}{\textbf{VRAI}}  & derivé : beneficie(ali, TemoignageFav.) \\
R3 & AccuseRecidiviste(omar)      & \textcolor{codegreen}{\textbf{VRAI}}  & derivé : a\_antecedent(omar, Condamnation) \\
R4 & CasComplexe(karim)           & \textcolor{codegreen}{\textbf{VRAI}}  & karim:AccuseMineur $\Rightarrow$ CasComplexe \\
R5 & CasSimple(ali)               & \textcolor{codegreen}{\textbf{VRAI}}  & asserté en ABox \\
R6 & instances(CasComplexe)       & \{omar, karim\}                       & omar(récidiviste) + karim(mineur) \\
R7 & instances(CasSimple)         & \{ali\}                               & asserté \\
R8 & AccuseRecidiviste $\sq$ Accuse & \textcolor{codegreen}{\textbf{VRAI}} & par définition de l'équivalence \\
R9 & Disjoint(CasSimple, CasComplexe) & \textcolor{codegreen}{\textbf{VRAI}} & $\dlcomp{R} \dlinter R = \bot$ \\
R10 & Disjoint(Juge, Procureur)   & \textcolor{red}{\textbf{FAUX}}        & tous deux Magistrats, non déclarés disjoints \\
\bottomrule
\end{tabularx}
\end{table}

\section{Implémentation avec OWL API 4.5}

\begin{lstlisting}[caption={Construction de la TBox et raisonnement HermiT}]
// Gestionnaire et ontologie
OWLOntologyManager mgr = OWLManager.createOWLOntologyManager();
OWLOntology onto = mgr.createOntology(IRI.create("http://rcr1/musique"));
OWLDataFactory f = mgr.getOWLDataFactory();

// Classes
OWLClass Artiste  = f.getOWLClass(IRI.create(NS + "Artiste"));
OWLClass Label    = f.getOWLClass(IRI.create(NS + "Label"));
OWLObjectProperty signeChez =
    f.getOWLObjectProperty(IRI.create(NS + "signe_chez"));

// TBox : ArtisteSigné ≡ Artiste ⊓ ∃signe_chez.Label
OWLClass ArtisteSigne = f.getOWLClass(IRI.create(NS + "ArtisteSigne"));
mgr.addAxiom(onto, f.getOWLEquivalentClassesAxiom(ArtisteSigne,
    f.getOWLObjectIntersectionOf(Artiste,
        f.getOWLObjectSomeValuesFrom(signeChez, Label))));

// ABox : signe_chez(soolking, universal)
OWLNamedIndividual soolking = f.getOWLNamedIndividual(IRI.create(NS+"soolking"));
OWLNamedIndividual universal = f.getOWLNamedIndividual(IRI.create(NS+"universal"));
mgr.addAxiom(onto, f.getOWLObjectPropertyAssertionAxiom(
    signeChez, soolking, universal));

// Raisonneur HermiT
OWLReasoner reasoner = new ReasonerFactory().createReasoner(onto);

// R6 : instances de ArtisteSigné
Set<OWLNamedIndividual> signes =
    reasoner.getInstances(ArtisteSigne, false).getFlattened();
// → {soolking, feu}

// R8 : ArtisteMajor ⊑ ArtisteSigné
boolean sub = reasoner.isEntailed(
    f.getOWLSubClassOfAxiom(ArtisteMajor, ArtisteSigne));
// → true
\end{lstlisting}
